\subsection{The precise separation consequences for the published theorems}

In this subsection $\Pi_Z^\Sigma$ always retains the entire arithmetic
quotient $G_k$. For a proper curve it is the geometrically pro-$\Sigma$
\'etale group; for an open curve it is the geometrically pro-$\Sigma$
tame group. The set $\Sigma$ may contain $p$, but must contain some
$r\ne p$. The relative comparison used in this scope is
\statementref{A.CF12}, not a claim that the full pro-$\Sigma$
configuration fibre sequence is exact.

\paragraph{The exceptional set in the source.}
Sa\"idi--Tamagawa 2017, Definitions~1.1, 1.2 and 1.4, first takes
the exceptional locus of the map from points of the universal
pro-cover to \emph{commensurability classes} of decomposition
subgroups, then takes its image $E_Z$ in the closed points of the
compactification. Its ``almost injective in the strong sense'' means
finite exceptional locus for a map of sets. In contrast,
Definition~2.9 calls a proper curve almost $\Sigma$-separated when
the residue fields of its exceptional points lie in a proper
algebraic subfield of $\bar k$. Proposition~2.12 assumes
$J_Z$-largeness and proves this latter support statement; it does
not assert that $E_Z$ is finite. The earlier cofinite result is
ST2009 Proposition~1.8(vi), with almost injectivity of the
closed-point decomposition map in Corollary~1.11(iii).
These are distinct assertions.\citep[Definitions~1.1, 1.2, 1.4,
2.6, 2.9; Proposition~2.12]{LA_ST2017}
\citep[Proposition~1.8(vi), Corollary~1.11]{LA_ST2009}

\begin{claim}{C.LIT5}{Proved: the refined proper theorem with every point labelled}{C1, C5}
Let $X/k_X$ and $Y/k_Y$ be proper smooth geometrically connected
curves of genus at least two over finite fields. Let $\Sigma_X,
\Sigma_Y$ each contain a prime different from the respective
characteristic, and suppose
\[
 \ker\!\left(G_{k_X}\longrightarrow
       \Aut\!\left(\prod_{s\notin\Sigma_X}T_sJ_X\right)\right)
 \ne1,
 \qquad
 \ker\!\left(G_{k_Y}\longrightarrow
       \Aut\!\left(\prod_{s\notin\Sigma_Y}T_sJ_Y\right)\right)
 \ne1.
\]
At the characteristic prime $T_p$ denotes the \'etale Tate module.
For every isomorphism $\alpha:\Pi_X^{\Sigma_X}\simeq
\Pi_Y^{\Sigma_Y}$, the two prime sets, characteristics and field
cardinalities coincide under the group-theoretic identifications,
and $\alpha$ preserves Frobenius. Moreover:
\begin{enumerate}
\item $E_X=E_Y=\varnothing$, and the same is true on every finite
connected \'etale covering arising from an open subgroup;
\item the decomposition-group construction gives a unique bijection
$\phi:X^{\cl}\simeq Y^{\cl}$, preserving residue-field degrees and
carrying $D_x$ to $D_{\phi(x)}$ up to conjugacy at every point;
\item every corresponding pair of finite marked sets $S,T$ may be
used in the cuspidally abelian reconstruction of ST2017
Corollary~3.3; the Green-compatibility assertion in Proposition~3.12
holds at all points and at every corresponding finite covering;
\item $\alpha$ is induced by the unique isomorphism of curves with
their corresponding pro-coverings in ST2017 Theorem~4.22.
\end{enumerate}
The fourth assertion was already the global conclusion of that
theorem. The improvement is the removal of the exceptional-point
restrictions in the preceding reconstruction statements.
\end{claim}
\begin{proof}
The displayed kernel conditions are exactly $J_X$-largeness and
$J_Y$-largeness in Definition~2.6. Proposition~1.7 recovers the
geometric subgroup and the prime sets. Lemma~2.7 gives largeness
relative to the constant fields, and Proposition~2.5 gives equality
of the characteristics and cardinalities and preservation of
Frobenius. Proposition~2.12 gives almost $\Sigma$-separation.
\statementref{A.CR10}, using the arbitrary-prime-set comparison
\statementref{A.CF12} and the retained boundary, upgrades it to
$E=\varnothing$. By ST2017 Proposition~2.11(i), an \'etale
$\Sigma$-cover has exceptional set equal to the inverse image of
the original exceptional set, so emptiness holds at every such
level, including constant extensions. No largeness hypothesis on
the new Jacobian is silently added.

Theorem~1.9 reconstructs the set of decomposition subgroups from
the Frobenius group. Emptiness makes the association with points
injective, hence gives $\phi$ on all closed points. The arithmetic
image of $D_x$ is $G_{k(x)}$, which proves the degree assertion.
Corollary~3.3 and Proposition~3.12 now have all their original
hypotheses with $X^{\cl}\setminus E_X=X^{\cl}$ and similarly
for $Y$ and for every covering. Substitution proves precisely the
marked-point and Green-compatibility assertions. Finally the
hypotheses of Theorem~4.22 are exactly those displayed here,
so its curve and pro-covering isomorphism applies.
\citep[Proposition~1.7, Theorem~1.9, Proposition~2.5,
Lemma~2.7, Propositions~2.11--2.12, Corollary~3.3,
Proposition~3.12, Theorem~4.22]{LA_ST2017}
\end{proof}

The $J$-largeness assumption includes every cofinite prime set and
some infinite-complement sets described in ST2017 Remark~2.8.
The present argument does not replace this condition by an arbitrary
nonempty $\Sigma$. In particular, the singleton-prime collisions
of \statementref{A.CR13} prevent such a replacement. Nor does
$E=\varnothing$ remove the omitted-prime torsion
$J(k)\{\Sigma'\}$ or the quotient by
$k^\times\{\Sigma'\}$ in the pseudo-function reconstruction
(Definition/Remark~4.1 and Theorem~4.2). The finite-index step in
the proof of Theorem~4.22 uses Remark~4.7's separate control of
complementary-prime torsion along a constant-field tower. Point
separation alone does not provide that control.

\begin{claim}{C.LIT6}{Proved: the exact affine hypothesis is retained}{C1, C5}
Let $U,V$ be hyperbolic curves over finite fields, with an
isomorphism of geometrically pro-$\Sigma$ tame groups. Assume,
as in ST2017 Theorem~4.23, that there are corresponding open
subgroups defining covers $U',V'$ whose smooth compactifications
$X',Y'$ have genus at least two, with $\Sigma$ respectively
$J_{X'}$-large and $J_{Y'}$-large. Then $E_U=E_V=\varnothing$,
all closed-point decomposition labels are unique, and the unique
scheme isomorphism with its pro-coverings in Theorem~4.23 has
no exceptional-point restriction. The matched-open-subgroup
hypothesis is part of this assertion.
\end{claim}
\begin{proof}
By C.LIT5, $E_{X'}=E_{Y'}=\varnothing$. If two interior points
of the universal pro-cover of $U$ had commensurate decomposition
groups, restrict to their common open arithmetic subgroup and to
the given open subgroup for $U'$. After a finite constant
extension they give distinct rational points of $U'$ with equal
section subgroups. They remain distinct on $X'$, and projection
to $\Pi_{X'}^\Sigma$ contradicts its all-extension separation.
Two distinct upstairs points above the same geometric point are
already excluded by ST2017 Proposition~1.3. Cusp decomposition
groups are recognized by their nontrivial inertia, and their
commensurability separation is the cusp assertion of ST2009
Proposition~1.8(i),(ii), recalled in ST2017 Theorem~1.8. Thus
no interior or cuspidal exception remains. The same argument
applies to $V$. Theorem~4.23 then applies with its original
matched-cover hypotheses. This proves both the new separation
clause and the already published geometric conclusion.
\citep[Proposition~1.3, Theorems~1.8 and 4.23]{LA_ST2017}
\end{proof}

\begin{claim}{C.LIT7}{Proved: the prime-to-characteristic theorem and its simplified restrictions}{C2, C5}
Let $U,V$ be hyperbolic curves over finite fields, and let
$\Sigma_U^\dagger$ and $\Sigma_V^\dagger$ consist of every prime
different from their respective characteristics. Every isomorphism
of their geometrically pro-$\Sigma$ tame groups arises from a unique
isomorphism of curves and corresponding pro-coverings. All their
closed-point decomposition maps are injective, including after
every constant extension and every open-subgroup covering. For
proper $X,Y$, the restrictions to complements of exceptional sets
in ST2009 Corollary~2.3, Theorem~3.6 and Proposition~3.8 can be
replaced by all points.
\end{claim}
\begin{proof}
For $\Sigma=p'$, \statementref{A.R1} gives proper point
separation after every finite extension and
\statementref{A.CR9} gives $E_X=\varnothing$. If $p$ is also
retained, a collision would survive in the prime-to-$p$ quotient,
so it is excluded as well. At every proper \'etale level the
curve still has genus at least two; the same proof applies there.
For a nonproper hyperbolic curve choose a corresponding finite
tame cover whose compactification has genus at least two, as in
the proof of ST2009 Corollary~3.10, and use the argument of
C.LIT6. The cusp separation is unchanged.

The exact published Isom statements are ST2009 Theorem~3.9
for proper curves and Corollary~3.10 for arbitrary hyperbolic
curves. They permit either all primes or all primes except $p$.
Their hypotheses are met here. In the proper reconstruction,
Corollary~2.3 now applies to arbitrary finite $S,T$;
$\mathcal O_{E_X}=K_X$ and $\mathcal O_{E_Y}=K_Y$ in
Theorem~3.6; and Proposition~3.8 has no excluded points at any
finite covering. The proof of Theorem~3.9 is the formal
combination of Theorem~3.7 with Proposition~3.8, and the proof
of Corollary~3.10 descends compatible proper isomorphisms from
the cofinal family of such covers. These are the precise places
where the separation result removes restrictions.
\citep[Corollary~2.3, Theorems~3.6--3.7, Proposition~3.8,
Theorem~3.9, Corollary~3.10]{LA_ST2009}

This simplification does not set $J_X(k_X)[p^\infty]$ equal
to zero. The $p$-power correction to functions in Theorem~3.6(ii)
therefore remains necessary. Similarly, the reconstruction of
additive structure is still a separate part of the proof; an
injective point-label map alone is not an Isom theorem.
\end{proof}

\paragraph{What becomes of the proposed configuration route.}
ST2009 arXiv Remark~3.13, journal Remark~3.14, would first obtain
an isomorphism of the \emph{full} prime-to-$p$ punctured groups
from the proposed full fibre exactness, then apply the published
Isom theorem to infer equality of the punctures. The full
exactness premise is false by Part~A. Replacing the group by its
cuspidally pro-$r$ quotient does not permit that last application
of the full-group Isom theorem. What does work is the distinct
chain
\[
 \text{relative exactness and diagonal transport}
 \Longrightarrow
 \text{the cuspidal criterion of Part A}
 \Longrightarrow E_X=\varnothing.
\]
This chain is independent of the published Isom conclusion.
Combining it with the explicitly retained cuspidally abelian
and fixed-$r$ reconstruction inputs gives C.LIT7 by the source's
function-reconstruction argument. It is a corrected route to
point separation followed by the known Isom proof, rather than
a restoration of the false full fibre premise.

\begin{claim}{C.LIT8}{Proved: specific Mochizuki separation hypotheses become automatic}{C5}
For proper hyperbolic curves over finite fields and
$\Sigma^\dagger=p'$, the $\Sigma$-separation assumption in
Mochizuki's original Definition~1.18, Proposition~2.6(ii),(iii)
and Corollary~3.11 is automatic. Thus a Frobenius-preserving
isomorphism is point-theoretic, arbitrary finite corresponding
marked sets are allowed in Theorem~2.5(i), the indicated proper
and cuspidally abelian cusp decomposition groups are
commensurably terminal, and the isomorphism is totally globally
Green-compatible. The quotient repairs of A.LIT2
are used at the fibre-identification inputs.
\end{claim}
\begin{proof}
Definition~1.18(i) asks for closed-point decomposition
injectivity on every connected finite \'etale covering arising
from an open subgroup. C.LIT7 proves exactly this requirement.
Decomposition-group recognition and this injectivity change
quasi-point-theoreticity into point-theoreticity under
Definition~1.18(ii), as in Remark~1.18.2(ii). Theorem~2.5(i)
assumes point-theoreticity and gives the marked-set assertion.
Proposition~2.6(ii),(iii) now have their separation assumption
satisfied and give the two commensurable terminality assertions.
Finally Corollary~3.11 assumes finite fields,
$\Sigma^\dagger=p'$ and separation; all three conditions hold.
Its fixed-prime fibre uses are in the repaired scope of A.LIT2.
These are applications of the stated source results with their
actual assumptions, not an assertion that every conclusion of
the entire paper has been reverified.
\citep[Definition~1.18, Remark~1.18.2, Theorem~2.5(i),
Proposition~2.6(ii),(iii), Corollary~3.11]{LA_Mochizuki2007}
\end{proof}

For reference, the original/LaTeX numbering pairs in the last
claim are: Definition~1.18/Definition~1.5;
Theorem~2.5/Theorem~2.1; Proposition~2.6/Proposition~2.3;
Corollary~3.11/Corollary~3.1. The source's
Theorem~2.5(ii) assumes $\Sigma$ is \emph{all} primes in order
to recover principal divisors exactly. Separation does not remove
that hypothesis. Its original Theorem~3.10 (LaTeX Theorem~3.1)
already permits arbitrary matching rational points; the new
fact makes the matching point unique, rather than supplying a
previously absent arbitrary-point existence assertion.

\paragraph{Bounds on the consequence claim.}
Wakabayashi's reconstruction theorems use a single invertible
prime. Their pro-$l$ point labels cannot be declared injective
from a theorem about the larger prime-to-$p$ group; Part~A gives
actual singleton-prime collisions. The virtual-curve construction
in Sun 2025 already applies to arbitrary sections, and Sun 2026's
specified open-injection criterion has its own geometricity
hypotheses. Point separation neither constructs a nongeometric
section nor realizes its virtual extension as a punctured
algebraic curve. Finally no arbitrary-homomorphism anabelian
theorem follows merely by replacing exceptional sets with the
empty set in the checked Isom statements. These limits are part
of the scope of C1, C2 and C5.

\subsection{The separation-hypothesis census in the remaining named papers}

\begin{claim}{C.LIT9}{Proved}{Exact removal of the named separation conditions}
In the original Mochizuki 2007 text, the explicit separation or
point-theoreticity conditions relevant to C5 are the loci in the
following table. For finite fields with $\Sigma^\dagger=p'$,
point-theoreticity and separation at these loci follow from
\statementref{A.R1} and its all-cover corollaries, as applied in
C.LIT7--8. Other hypotheses in the same statements remain.
Wakabayashi 2016 and the two Sun preprints have no additional
unproved finite-field prime-to-$p$ separation condition whose
removal strengthens their stated theorems. The distinctions in
the table are part of the conclusion.
\end{claim}
\begin{center}\small
\begin{longtable}{p{0.25\linewidth}p{0.33\linewidth}p{0.34\linewidth}}
\toprule
Locus&Condition or use&Consequence and retained restriction\\\midrule
Mochizuki Definition~1.18; Remark~1.18.2(i)--(iii)&
Separation on every finite open-subgroup cover; passage from
quasi-point-theoretic to point-theoretic; possible failure of
maximality or terminality of point decomposition groups.&
Separation and unique point labels hold at every such level.
For a Frobenius-preserving isomorphism the matching point is
unique. The terminality conclusion is Proposition~2.6 below.\\
Proposition~2.3(ii)&
Proves separation when $\Sigma$ is all primes.&
Its finite-field point-separation conclusion now also holds for
$\Sigma=p'$. Its assertion about detecting principal divisors
in (i) is a different statement and is not extended by this
substitution.\\
Definition~2.4; Theorem~2.5(i)&
Point-theoreticity supplies matching point sets used in Green
compatibility and marked cuspidally abelian reconstruction.&
All finite corresponding sets are allowed, functorially, with
the same cusp and proper-quotient data. The auxiliary definition
of Green compatibility is unchanged.\\
Theorem~2.5(ii)&
Adds $\Sigma=\mathrm{Primes}$ to reconstruct principal
divisors and the stated push-forward of multiplicative groups.&
Point-theoreticity is automatic in the finite-field setting,
but the all-prime hypothesis is retained. Separation alone does
not remove complementary-primary torsion.\\
Proposition~2.6(ii),(iii)&
Assumes $\Sigma$-separation for commensurable terminality in
the proper group and in the cuspidally abelian punctured group.&
Both terminality conclusions hold without an additional
separation assumption in the finite-field prime-to-$p$ cases.\\
Remark~2.6.1; Corollary~2.7&
The remark compares an already given punctured-group
isomorphism with the reconstructed abelian one. The corollary
assumes all primes, point-theoreticity, and total local principal
Green compatibility.&
The point-label condition is automatic as above. Existence of
the punctured-group isomorphism in the remark is not deduced
from separation. In the finite-field all-prime corollary,
Corollary~3.11 supplies Green compatibility as in C.LIT8;
the Isom conclusion was already part of the source's program.\\
Corollary~2.8; Proposition~2.10(iii)&
The former is a point-theoretic affine Isom theorem with all
primes; the latter supplies point-theoreticity for quasi-Belyi
curves over nonarchimedean local fields.&
The finite-field all-prime premise of the former is automatic;
the prime-to-$p$ version is obtained by the separately checked
ST2009 argument, C.LIT7. The local-field result has no new
finite-field consequence from R1.\\
Theorem~3.10; Corollary~3.11&
The theorem assumes matching rational point decomposition
groups, not separation. The corollary assumes finite fields,
$\Sigma^\dagger=p'$, and separation.&
The matching point in the theorem becomes unique; arbitrary
matching-point existence was already present. The corollary's
separation condition disappears and total global Green
compatibility follows.\\
Wakabayashi Theorems~3.9,4.4; Corollary~3.10 and
Remark~3.10.1&
One invertible prime $l$; matching decomposition groups in the
punctured theorem; the remark explicitly allows several
matching closed points.&
The source does not assume a mixed-prime separation condition.
R1 cannot make its single-prime point correspondence unique;
the collisions in Part~A prohibit that inference.\\
Sun 2025, Definition~3.1, Proposition~3.7, Corollary~3.8,
Theorems~A--C&
Virtual pullbacks are defined for arbitrary sections in
characteristic zero. Geometric sections identify their actual
fibres; the main reconstruction statements use number fields
or mixed-characteristic local fields.&
There is no conditional finite-field separation clause to
remove. Point separation does not prove that a virtual curve
has an algebraic realization or that its section is geometric.\\
Sun 2026, Lemma~5.6; Theorem~5.7 and Corollary~5.8&
Lemma~5.6(i) uses the already known injectivity of points into
full sections over an arithmetic field; (ii) uses an additional
open-immersion geometricity theorem. The later results impose
specific virtual covering and cusp conditions.&
No such condition is supplied merely by R1. The source's
arithmetic fields are number fields or mixed-characteristic
local fields, and its groups are full \'etale groups.
None of these assertions becomes a finite-field prime-to-$p$
section-surjectivity theorem.\\
\bottomrule
\end{longtable}
\end{center}
\begin{proof}
We searched the complete opened texts for the source's
separation, point-theoreticity, point-injectivity and geometricity
terms, then read the surrounding numbered statements and their
proof uses. In Mochizuki's original text, the explicit
$\Sigma$-separation occurrences are Definition~1.18,
Remark~1.18.2, Proposition~2.3(ii), Proposition~2.6(ii),(iii)
and Corollary~3.11. The point-theoreticity occurrences add
Definition~2.4, Theorem~2.5, Remark~2.6.1, Corollaries~2.7--2.8
and Proposition~2.10(iii). Theorem~3.10 requires corresponding
point groups, but not separation. This produces the displayed
finite list; no unexamined whole-paper assertion is substituted
for it.

C.LIT7 proves separation on every required cover, C.LIT8 proves
the resulting point-theoreticity and terminality, and A.LIT2
and A.LIT8 supply the repaired quotient fibre inputs. Substituting
these established assertions into the displayed hypotheses gives
the consequences in the first eight rows. In particular, the
proof of Theorem~2.5(i) restricts the reconstructed cuspidally
abelian configuration extension to the matched point groups;
the repaired abelian fibre is exactly what this restriction
requires. The all-prime requirement in (ii) is separate and
has not been removed. The proof of Corollary~3.11 uses the
fixed-$l$ reconstruction for each $l\in\Sigma^\dagger$
and recognizes noncuspidal decomposition groups there. A.LIT8
repairs that input, while all-cover separation makes each of
the matching labels unique. This is precisely total global
Green compatibility.

Wakabayashi's Remark~3.10.1 expressly distinguishes existence
of a matching point from uniqueness. Its surrounding theorems
remain in the single-prime category described in A.LIT4, so
the larger prime-to-$p$ result does not alter their hypotheses.
Sun 2025 constructs the virtual pullback before making any
assumption that its section is a point. Sun 2026 defines an
arithmetic field to mean a number field or a mixed-characteristic
local field; its Lemma~5.6(i) invokes the injectivity portion
of the section conjecture in that already known full-group
setting. Its part (ii) and later geometricity criteria have
additional inputs, not a missing point-separation premise.
These readings prove the last three rows and complete the
stated separation-condition census. They also explain why no
new arbitrary-homomorphism theorem or section-surjectivity
theorem has been inferred here.
\citep[Definition~1.18, \S2, Theorem~3.10,
Corollary~3.11]{LA_Mochizuki2007}
\citep[Theorems~3.9,4.4, Remark~3.10.1]{LA_Wakabayashi2016}
\citep[Definition~3.1, Proposition~3.7, Corollary~3.8,
Theorems~A--C]{LA_Sun2025}
\citep[\S2, Lemma~5.6, Theorem~5.7, Corollary~5.8]{LA_Sun2026}
\end{proof}
